% ==============================================================================
% T0/FFGFT Theory: Document 230
% Title: Translatability FFGFT ↔ Hilbert-space quantum mechanics
% Author: Johann Pascher
% Date: May 2026
% References: Documents 035 (QM foundations), 147 (Quantum Computing),
%             174 (Spin qubit), 175 (Qubit state spaces), 202 (Field theory)
% ==============================================================================

\section*{Preamble}

This document formulates the \textbf{complete translatability} between
the FFGFT mode formalism and standard Hilbert-space quantum mechanics.
The translation is not an approximation or specialisation, but a
bijective structural statement: every QM calculation can be mapped
into FFGFT language, and conversely every FFGFT statement about
quantum systems is reproducible in Hilbert-space language.

\textbf{This is methodologically important:} it means that FFGFT does
not \emph{replace} quantum mechanics, but \emph{extends} it. QM is a
subset of FFGFT (see Doc.~202 §22). Practitioners familiar with
Hilbert-space language can use FFGFT without changing their tools --
the bridge is the translation table below.

% ==============================================================================
\section*{Where we are, geometrically}
% ==============================================================================

Before the formal translation tables, a brief narrative picture of
the underlying geometry. This matters because we are about to work
with \emph{cylindrical coordinates} -- and the cylinder is not the
full FFGFT geometry, but a convenient reduction.

\paragraph{The full FFGFT geometry is four-dimensional.}
The fundamental space on which all FFGFT statements live is the
4-torus $T^4 = T^3 \times S^1$ -- three spatial winding directions
and one temporal winding (Doc.~202 §13, Doc.~210). The wave
function of a mode $\varphi_{n,l,j}$ is a function on this 4D
torus. From this geometry follow all masses, the fine-structure
constant, cosmology -- everything that distinguishes FFGFT from
pure quantum mechanics.

\paragraph{For a single qubit, less suffices.}
When we want to describe just a single qubit -- a two-level system
at fixed time, without mass dynamics -- most of these winding
directions are not active. Two rotational degrees of freedom remain:
one encoding \enquote{where between $|0\rangle$ and $|1\rangle$ we
are} (this becomes our $z$), and one for the phase (this becomes
our $\theta$). The natural space for these two rotational degrees
of freedom is a 2-torus $T^2 = S^1 \times S^1$.

\paragraph{From 2-torus to cylinder.}
If we now assume that one major radius $R$ of the 2-torus is very
large compared to the tube radius $r$, then the local geometry is
no longer curved -- it looks like a cylinder. The cylinder
$\mathbb{R} \times S^1$ is thus the limit $R \to \infty$ of the
2-torus. Topologically one loop is \enquote{cut open}: from two
circles we get a line plus a circle. In practice this means: in
cylindrical coordinates $(z, r, \theta)$, $z$ is a \emph{linear}
axis (from $-1$ to $+1$), no longer cyclic.

\paragraph{From cylinder to Bloch sphere.}
If we contract the cylinder top and bottom to a point -- the two
\enquote{poles} -- we obtain a sphere $S^2$. This is the famous
\textbf{Bloch sphere} of standard quantum mechanics. It is the
second, further reduction: cylinder with merged ends.

\paragraph{Hierarchy of reductions.}
\begin{center}
\renewcommand{\arraystretch}{1.4}
{\small
\begin{tabular}{>{\raggedright\arraybackslash}p{3.0cm}>{\raggedright\arraybackslash}p{1.4cm}>{\raggedright\arraybackslash}p{4.0cm}}
\toprule
\textbf{Geometry} & \textbf{Dim.} & \textbf{Use} \\
\midrule
4-torus $T^4$ & 4 & full FFGFT (masses, $\alpha$, cosmology) \\
3-torus $T^3$ & 3 & mode space at fixed time \\
2-torus $T^2$ & 2 & single qubit, without reduction \\
Cylinder $\mathbb{R} \times S^1$ & 2 & limit $R \to \infty$, local approximation \\
Bloch sphere $S^2$ & 2 & poles contracted, standard QM \\
\bottomrule
\end{tabular}
}
\renewcommand{\arraystretch}{1.0}
\end{center}

\paragraph{What the reductions cost.}
Each step of reduction loses something. From the 4-torus to the
3-torus one loses the time winding -- and with it all mass
statements (Doc.~210). From the 3-torus to the 2-torus one loses a
spatial winding -- and with it the connection to the third lepton
generation and many spin structures. From the 2-torus to the
cylinder one loses \emph{a compact loop} -- this is the weakest
reduction, and exactly here sits the \textbf{fractal correction
factor} $K_{\text{frak}} \approx 0.9867$ that appears in the
bit-flip rotations of Doc.~175. $K_{\text{frak}}$ is the
\enquote{memory} of the cylinder of its toroidal origin.

\paragraph{What we translate here.}
The translation tables below apply to the \emph{qubit sector},
i.e.\ the lower three rows of the hierarchy -- 2-torus, cylinder,
Bloch sphere. Standard quantum mechanics works on the Bloch
sphere; FFGFT works on the cylinder (or the 2-torus, if one takes
$K_{\text{frak}}$ seriously). The translation between these levels
is lossless, provided one carries the small geometric corrections
($K_{\text{frak}}$, $\Delta\text{CHSH}$) along.

The full FFGFT geometry -- the 4-torus -- lies one step higher and
provides what Hilbert-space QM cannot provide: the values of the
input parameters (masses, $\alpha$, $\xi$). This is discussed in
§10 below.

% ==============================================================================
\section{\texorpdfstring{States: $(z, r, \theta)$ vs.\ $(\alpha, \beta)$}{States: (z,r,theta) vs. (alpha,beta)}}
% ==============================================================================

\subsection{Hilbert-space representation}

In standard QM a qubit is a vector in $\mathbb{C}^2$:
\begin{equation}
  |\psi\rangle = \alpha|0\rangle + \beta|1\rangle,
  \qquad
  |\alpha|^2 + |\beta|^2 = 1,
  \qquad
  \alpha, \beta \in \mathbb{C}.
\end{equation}
With global phase freedom this is a point on the Bloch sphere
(Hopf fibration $S^3 \to S^2$).

\subsection{FFGFT representation}

In FFGFT the same qubit is a point $(z, r, \theta)$ on the
cylinder/torus:
\begin{equation}
  z \in [-1, 1], \quad r \in [0, 1], \quad \theta \in [0, 2\pi),
  \qquad
  z^2 + r^2 = 1.
\end{equation}
$z$ is the projection onto the computational basis axis, $r$ the
radial superposition amplitude, $\theta$ the phase (Doc.~147 §1.2).

\subsection{Bijective translation}

The two representations are linked by:
\begin{equation}
  \boxed{
  \alpha = \sqrt{\tfrac{1+z}{2}} \cdot e^{i\theta/2},
  \qquad
  \beta = \sqrt{\tfrac{1-z}{2}} \cdot e^{-i\theta/2}.}
  \label{eq:state-bridge}
\end{equation}
Conversely:
\begin{equation}
  z = |\alpha|^2 - |\beta|^2,
  \qquad
  r = 2|\alpha\beta|,
  \qquad
  \theta = \arg(\alpha) - \arg(\beta).
\end{equation}

\textbf{Consistency checks:}
\begin{itemize}[leftmargin=2em,itemsep=2pt,topsep=2pt]
  \item $|\alpha|^2 + |\beta|^2 = \tfrac{1+z}{2} + \tfrac{1-z}{2} = 1$. \checkmark
  \item $z^2 + r^2 = (|\alpha|^2 - |\beta|^2)^2 + 4|\alpha|^2|\beta|^2 = 1$. \checkmark
  \item $|0\rangle$ corresponds to $(z,r,\theta) = (1, 0, *)$. \checkmark
  \item $|1\rangle$ corresponds to $(z,r,\theta) = (-1, 0, *)$. \checkmark
\end{itemize}

\textbf{Conceptual difference:} in QM, $|\alpha|^2$ is a
\emph{probability} (Born rule). In FFGFT, $r^2 = 1 - z^2$ is an
\emph{energy density} of the radial configuration. The two
quantities are \emph{numerically equal} at large statistics
(Doc.~147), but conceptually different: FFGFT is a deterministic
geometric model whose statistical predictions coincide with the
Born rule.

% ==============================================================================
\section{Operators: Pauli matrices vs.\ rotations}
% ==============================================================================

\subsection{Hilbert-space operators}

The three Pauli matrices generate all single-qubit operations:
\begin{equation}
  \sigma_x = \begin{pmatrix}0 & 1 \\ 1 & 0\end{pmatrix},
  \quad
  \sigma_y = \begin{pmatrix}0 & -i \\ i & 0\end{pmatrix},
  \quad
  \sigma_z = \begin{pmatrix}1 & 0 \\ 0 & -1\end{pmatrix}.
\end{equation}
Algebra: $\sigma_i \sigma_j = \delta_{ij} \mathbb{1} + i\epsilon_{ijk}\sigma_k$.

\subsection{FFGFT rotations}

On the Bloch cylinder, the three Pauli matrices correspond to three
geometric rotations (Doc.~175 §3):
\begin{itemize}[leftmargin=2em,itemsep=2pt,topsep=2pt]
  \item $\sigma_z$ -- rotation about the cylinder axis (phase rotation $\theta \to \theta + \pi$).
  \item $\sigma_x$ -- 2D rotation in the $(z, r)$ plane by angle $\pi$ (bit flip).
  \item $\sigma_y$ -- orthogonal 2D rotation, combining $\sigma_x$ and phase rotation.
\end{itemize}

\subsection{Translation table}

\begin{center}
\renewcommand{\arraystretch}{1.4}
{\small
\begin{tabular}{>{\raggedright\arraybackslash}p{1.6cm}>{\raggedright\arraybackslash}p{3.2cm}>{\raggedright\arraybackslash}p{3.6cm}}
\toprule
\textbf{Operator} & \textbf{Hilbert-space matrix} & \textbf{FFGFT geometry} \\
\midrule
$\sigma_z$ & $\begin{pmatrix}1 & 0 \\ 0 & -1\end{pmatrix}$
           & Phase rotation $\theta \to \theta + \pi$ \\
$\sigma_x$ & $\begin{pmatrix}0 & 1 \\ 1 & 0\end{pmatrix}$
           & $(z, r)$-rotation by $\pi K_{\text{frak}}$ \\
$\sigma_y$ & $\begin{pmatrix}0 & -i \\ i & 0\end{pmatrix}$
           & Rotation with phase tilt \\
$H$ & $\frac{1}{\sqrt{2}}\begin{pmatrix}1 & 1 \\ 1 & -1\end{pmatrix}$
    & Swap $z \leftrightarrow r$, $\theta \to \theta + \pi/2$ \\
$T$ & $\begin{pmatrix}1 & 0 \\ 0 & e^{i\pi/4}\end{pmatrix}$
    & Phase rotation $\theta \to \theta + \pi/4$ \\
\bottomrule
\end{tabular}
}
\renewcommand{\arraystretch}{1.0}
\end{center}

\textbf{The only FFGFT-specific deviation} is the factor
$K_{\text{frak}} \approx 0.9867$ in the $\sigma_x$ rotation
(Doc.~175). This deviation follows from the fractal spacetime
dimension $D_f = 2.999867$ and is a \emph{testable prediction}: all
$\pi$-gates deviate systematically by $\Delta\alpha \approx 0.013\,\%$
from the ideal value. In Hilbert-space formalism this would be an
ad hoc correction modeled as hardware noise; in FFGFT it follows
from the geometry.

% ==============================================================================
\section{Quantum gates: complete translatability}
% ==============================================================================

\subsection{Single-qubit gates}

Every unitary single-qubit gate $U \in U(2)$ can be represented as a
Bloch-sphere rotation:
\begin{equation}
  U = e^{i\alpha} R_{\hat n}(\phi)
    = e^{i\alpha}(\cos(\phi/2)\mathbb{1} - i\sin(\phi/2)\, \hat n \cdot \vec\sigma).
\end{equation}
In FFGFT this is a rotation of the $(z, r, \theta)$ point about axis
$\hat n$ by angle $\phi$. The translation is exact: $U |\psi_\text{QM}\rangle$
corresponds to the rotation of the FFGFT cylinder point about $\hat n$
by $\phi$.

\subsection{Two-qubit gates}

The CNOT gate:
\begin{equation}
  \text{CNOT} = \begin{pmatrix}
    1 & 0 & 0 & 0 \\
    0 & 1 & 0 & 0 \\
    0 & 0 & 0 & 1 \\
    0 & 0 & 1 & 0
  \end{pmatrix}
\end{equation}
translates in FFGFT as: \emph{if $z_A = -1$ (control qubit on
$|1\rangle$), perform a bit flip on qubit B}. This is geometric
conditioning: the rotation of one cylinder point depends on the
$z$ value of the other. Topologically this corresponds to a
linkage of the two cylinders (entanglement).

\subsection{Universality}

The set $\{H, T, \text{CNOT}\}$ is universal for QM (Solovay-Kitaev
theorem). Since each of these gates is representable as a geometric
operation in FFGFT, FFGFT is \emph{fully universal for quantum
computation}. There is no QM statement that cannot be formulated in
FFGFT.

% ==============================================================================
\section{Tensor products and many-qubit spaces}
% ==============================================================================

\subsection{Hilbert-space construction}

For $n$ qubits the Hilbert space is
$\mathcal{H}_n = \bigotimes_{i=1}^{n} \mathbb{C}^2$ with dimension
$2^n$. A general state is:
\begin{equation}
  |\Psi\rangle = \sum_{x \in \{0,1\}^n} \alpha_x |x\rangle,
  \qquad
  \sum_x |\alpha_x|^2 = 1.
\end{equation}

\subsection{FFGFT product space}

In FFGFT $n$ qubits correspond to $n$ independent mode configurations
$(z_i, r_i, \theta_i)$ on a shared torus geometry. Entanglement is
represented as a \emph{topological linkage}: entangled qubits are
nodes of the mode configuration that cannot be reduced to product
form (Doc.~175 §4).

\subsection{Bell states}

The maximally entangled Bell state
\begin{equation}
  |\Phi^+\rangle = \tfrac{1}{\sqrt{2}}(|00\rangle + |11\rangle)
\end{equation}
has in FFGFT a geometric characterisation as a \emph{topologically
linked cylinder configuration}, in which the individual
$(z_i, r_i, \theta_i)$ values are not independent.

\textbf{CHSH prediction:} standard QM gives
$|\langle\text{CHSH}\rangle| \leq 2\sqrt{2}$ (Tsirelson bound).
FFGFT predicts a small, geometrically motivated deviation:
\begin{equation}
  \Delta\text{CHSH} \approx 10^{-5}
  \qquad\text{(geometric, from $\xi$ and $K_{\text{frak}}$)}.
\end{equation}
This deviation lies below the current hardware noise floor and is
an open experimental task for high-precision Bell tests
(Doc.~175 §6, Doc.~023).

\subsection{Exponential growth}

Both formalisms reproduce the exponential growth $\dim = 2^n$. In
Hilbert space through tensor product; in FFGFT through $n$
independent mode configurations with topological entanglements. The
number of real parameters per state is in both cases $2 \cdot 2^n - 2$.

% ==============================================================================
\section{Measurement}
% ==============================================================================

\subsection{Hilbert-space Born rule}

A measurement in the computational basis yields outcome $0$ with
probability $|\alpha|^2$ and outcome $1$ with probability $|\beta|^2$.
The state collapses to $|0\rangle$ or $|1\rangle$ respectively.

\subsection{FFGFT measurement as polar transition}

In FFGFT, measurement is \emph{not} a probabilistic collapse, but a
\emph{geometric interaction} that drives the $(z, r, \theta)$ point
to the nearest stable extremum ($z = +1$ or $z = -1$). This is a
deterministic motion in mode space -- but due to the unknown phase
$\theta$ and tiny variations in initial conditions, it is
\emph{statistically} indistinguishable from Born-rule behaviour.

\subsection{Numerical equivalence}

At large statistics (many repetitions), the FFGFT measurement
distribution converges to the Born rule:
\begin{equation}
  P_\text{FFGFT}(z = +1) = \tfrac{1+z}{2} = |\alpha|^2 = P_\text{QM}(0).
\end{equation}
The two formalisms yield \emph{identical} statistical predictions.

\textbf{Conceptual difference:} the Bell-test realism debate
resolves geometrically in FFGFT, since entanglement is represented
as a topological link -- not as nonlocal correlation. The
experimental predictions are the same.

% ==============================================================================
\section{Unitary evolution: Schrödinger equation}
% ==============================================================================

\subsection{Hilbert-space Schrödinger}

The time evolution in QM is:
\begin{equation}
  i\hbar \frac{\partial |\psi\rangle}{\partial t} = H |\psi\rangle.
\end{equation}

\subsection{FFGFT mode evolution}

In FFGFT the time evolution is the deterministic motion of the
mode configuration $\varphi_{n,l,j}$ on the torus
$T^3 \times \mathbb{R}$ (Doc.~202 §13). The corresponding
equation of motion -- the FFGFT mode equation -- reduces in the
low-energy limit ($E \ll E_\xi$) to the Schrödinger equation:
\begin{equation}
  i\hbar \frac{\partial \varphi_{n,l,j}}{\partial t}
  \;\xrightarrow{E \ll E_\xi}\;
  -\frac{\hbar^2}{2m}\nabla^2 \varphi_{n,l,j} + V \varphi_{n,l,j}.
\end{equation}
At higher energies, $\xi$ corrections appear, which in the
Hilbert-space formalism manifest as Hamiltonian modifications.

% ==============================================================================
\section{What is fully translatable}
% ==============================================================================

\textbf{All Hilbert-space statements map into FFGFT:}
\begin{itemize}[leftmargin=2em,itemsep=2pt,topsep=2pt]
  \item States, operators, gates (with or without $K_{\text{frak}}$ correction),
  \item Tensor products, entanglement, Bell states, GHZ states,
  \item Measurement (statistical predictions identical to Born rule),
  \item Unitary evolution (Schrödinger in low-energy limit),
  \item Algorithms (Shor, Grover, VQE, etc.).
\end{itemize}

\textbf{All FFGFT QM statements map into Hilbert space:}
\begin{itemize}[leftmargin=2em,itemsep=2pt,topsep=2pt]
  \item The $(z, r, \theta)$ configuration corresponds uniquely to
        the $(\alpha, \beta)$ vector (bridge (\ref{eq:state-bridge})),
  \item Geometric rotations correspond to unitary operators,
  \item Topological entanglements correspond to entangled
        Hilbert-space vectors.
\end{itemize}

\textbf{Consequence:} every computation that can be performed in one
formalism can be performed in the other -- with identical end
result. Practitioners can choose the formalism familiar to them.

% ==============================================================================
\section{What is FFGFT-specific and not representable in pure Hilbert space}
% ==============================================================================

The translatability holds for \emph{statements about quantum systems}.
It does not hold for \emph{statements about the input values} of
quantum mechanics:

\begin{itemize}[leftmargin=2em,itemsep=2pt,topsep=2pt]
  \item The geometric parameter $\xi_0 = 4/30000$ appears in the
        Hilbert-space formalism only as an external input. In FFGFT
        it follows from Higgs matching $\xi_0 = \lambda_h^2 v^2/
        (16\pi^3 m_h^2)$ (Doc.~006).
  \item The fermion masses $m_e, m_\mu, m_\tau, m_u, \ldots$ appear
        in Hilbert space as external Yukawa couplings. In FFGFT they
        all follow parameter-free from $m_i = r_i \xi_0^{p_i} v$
        (Doc.~006).
  \item The fine-structure constant $\alpha = \xi_0 \cdot E_0^2$ is
        derivable in FFGFT, an input in Hilbert space.
  \item The fractal correction factor $K_{\text{frak}} = 1 -
        100\xi_0 \approx 0.9867$ appears in FFGFT gates as a
        geometric deviation. In Hilbert space it would have to be
        modeled as hardware noise.
\end{itemize}

\textbf{This is the structural asymmetry:}
\begin{itemize}[leftmargin=2em,itemsep=2pt,topsep=2pt]
  \item Hilbert space: has all tools for QM calculations, but no
        explanation for the input values (masses, couplings,
        $\alpha$).
  \item FFGFT: provides both the tools for QM calculations
        \emph{and} the input values from a single geometry.
\end{itemize}

In this sense QM is a subset of FFGFT (see Doc.~202 §22): QM is
FFGFT, restricted to the question \emph{how systems behave}, without
the question \emph{why the constants take these values}.

% ==============================================================================
\section{Consequence: extension, not replacement}
% ==============================================================================

FFGFT does not replace quantum mechanics. It extends it to a level
where the input values themselves follow from geometry.

\textbf{For practitioners this means:}
\begin{itemize}[leftmargin=2em,itemsep=2pt,topsep=2pt]
  \item Those who calculate with Hilbert-space tools can continue to
        do so. The translation (\ref{eq:state-bridge}) and the
        operator table allow switching to the FFGFT formalism at any
        time, without loss.
  \item Those who calculate with the FFGFT mode formulation
        additionally have access to masses, couplings, and $\alpha$
        from the same geometry.
  \item Both languages are \emph{interoperable}. This is precisely
        the \enquote{translatability without contradiction} from
        Doc.~202~§22: other frameworks can map FFGFT statements into
        their language, provided they use the translation table.
\end{itemize}

\textbf{Methodological clarification:} FFGFT does not claim to be
the only correct representation of quantum mechanics. It is
\emph{one} representation that, in addition to the standard QM
statements, provides the values of the fundamental constants. Those
who do not need this added value can ignore FFGFT -- the standard
QM statements remain valid.

% ==============================================================================
\section{Summary}
% ==============================================================================

The translation table FFGFT $\leftrightarrow$ Hilbert-space QM is
complete:

\begin{enumerate}[leftmargin=2em,itemsep=2pt,topsep=2pt]
  \item \textbf{States:} bijective bridge
        $\alpha = \sqrt{(1+z)/2}\,e^{i\theta/2}$,
        $\beta = \sqrt{(1-z)/2}\,e^{-i\theta/2}$.
  \item \textbf{Operators:} Pauli matrices $\leftrightarrow$ Bloch
        rotations, with FFGFT correction $K_{\text{frak}}$ in
        $\sigma_x$.
  \item \textbf{Gates:} all standard gates ($H$, $T$, CNOT, ...) in
        FFGFT as geometric transformations.
  \item \textbf{Tensor products:} $2^n$ growth reproduced,
        entanglement as topological linkage.
  \item \textbf{Measurement:} statistically identical to Born rule,
        conceptually deterministic geometry.
  \item \textbf{Evolution:} Schrödinger equation as low-energy limit
        of FFGFT mode evolution.
\end{enumerate}

\textbf{FFGFT added value beyond Hilbert space:}
\begin{itemize}[leftmargin=2em,itemsep=2pt,topsep=2pt]
  \item parameter-free derivation of $\alpha$, $m_i$, $\xi$ from a
        single geometry,
  \item geometric explanation of the fractal correction factor
        $K_{\text{frak}}$ in $\pi$-gates (testable in
        high-precision quantum gate experiments),
  \item geometric resolution of the Bell-test realism debate
        (entanglement as topology, not as non-locality),
  \item small, geometrically motivated CHSH deviation $\sim 10^{-5}$
        as a testable prediction.
\end{itemize}

\textbf{Methodological position:} FFGFT confirms quantum mechanics
and extends it. It refutes nothing, but adds structural explanations
for the input values. The translatability $\leftrightarrow$ Hilbert
space is a \emph{consistency proof}: any internally consistent
theory should be expressible in multiple mathematical languages, and
the existence of such translations is evidence of structural
soundness.
